\subsubsection{Abrazadera deslizante con resorte y varilla rígida}

Una abrazadera de masa $m$ se desliza sin fricción a lo largo de una varilla rígida y sin masa, como se muestra en la \textbf{Figura \ref{fig:abrazadera_resorte}}. La abrazadera está unida a un resorte lineal con constante elástica $K$ y longitud natural $L$. Suponiendo ausencia de gravedad, determine la ecuación diferencial del movimiento de la abrazadera. Utilice los conceptos de trabajo y energía.

\begin{enumerate}
    \item Defina los marcos de referencia y sistemas de coordenadas adecuados.
    \item Analice la cinemática de la partícula en función de la coordenada $x$.
    \item Analice la cinética de la partícula utilizando el método de energía para obtener la ecuación diferencial de movimiento.
\end{enumerate}

\begin{figure}[!ht]
    \centering
    \begin{tikzpicture}[>=Latex, scale=1.5]
        % --- Parámetros ---
        \def\L{2.0}      % Distancia vertical L
        \def\xpos{1.2}   % Posición x de la masa
        \def\Wvarilla{2.5} % Media longitud de la varilla

        % --- 1. Soportes y Paredes ---
        % Soporte superior O
        \begin{scope}[shift={(0,\L)}]
            \fill[pattern=north east lines] (-0.4,0) rectangle (0.4,0.2);
            \draw[thick] (-0.4,0) -- (0.4,0);
            %\filldraw[fill=gris1!20] (0,0.1) arc (90:270:0.15) -- (0,-0.15) arc (-90:90:0.15);
            \filldraw (0,0) circle (1.5pt) node[below left, xshift=1pt] {$O$};
        \end{scope}

        % --- 2. Varilla Horizontal ---
        \draw[line width=7.4pt, black] (-\Wvarilla, 0) -- (\Wvarilla, 0);
        \draw[line width=6pt, gris1!30] (-\Wvarilla, 0) -- (\Wvarilla, 0);
        % Paredes laterales de la varilla
        \fill[pattern=north east lines] (-\Wvarilla-0.2, -0.4) rectangle (-\Wvarilla, 0.4);
        \draw[thick] (-\Wvarilla, -0.4) -- (-\Wvarilla, 0.4);
        \fill[pattern=north east lines] (\Wvarilla, -0.4) rectangle (\Wvarilla+0.2, 0.4);
        \draw[thick] (\Wvarilla, -0.4) -- (\Wvarilla, 0.4);


        % --- 3. Líneas de Referencia y Cotas ---
        \draw[dashed] (0,\L) -- (0, -0.5);
        \draw[<->] (-0.3, 0.1) -- (-0.3, \L) node[midway, left] {$L$};
        \draw[|->, thick] (0, -0.3) -- (\xpos, -0.3) node[midway, below] {$x$};

        % --- 4. Resorte ---
        \draw[decoration={aspect=0.3, segment length=3pt, amplitude=3pt, coil}, decorate, thick] 
            (0,\L) -- (\xpos, 0);
        \node[right, xshift=5pt] at (\xpos/2, \L/2) {$K$};

        % --- 5. Abrazadera (Masa m) ---
        \shadedraw[ball color=gris1!30] (\xpos, 0) ellipse (0.15 and 0.15);
        \node[above right, yshift=2pt] at (\xpos, 0) {$m$};

        % --- 6. Vector Gravedad (Ausencia indicada) ---
        %\draw[->, thick, myred] (2.2, 1.8) -- (2.2, 1.2) node[midway, right] {$\mathbf{g} = 0$};

    \end{tikzpicture}
    \caption{Sistema de una abrazadera de masa $m$ bajo la acción de un resorte lineal en una guía horizontal.}
    \label{fig:abrazadera_resorte}
\end{figure}